problem:  
Let \( (M^2,g) \) be a compact Riemann surface and \( (N,h) \) a compact Riemannian manifold.  Consider a sequence of approximate harmonic maps  
\[
u_k:M\to N,\qquad \tau(u_k)=f_k,
\]
with uniformly bounded energy \(E(u_k)\le C\) and \(f_k\to0\) weakly in \(H^{-1}(M,TN)\) (using an isometric embedding of \(N\) in \(\mathbb{R}^m\)).  
Suppose \(u_k\rightharpoonup u_\infty\) in \(H^1(M,N)\) and bubbling occurs at finitely many points \(x_1,\ldots,x_m\in M\).  
After rescaling at \(x_i\), assume one obtains bubbles \(\phi_{i,j}:S^2\to N\), each a harmonic sphere.

Denote by \(L_{\phi_{i,j}}\) the Jacobi operator on \(\phi_{i,j}^*TN\), and by
\[
\mathcal K_{\phi_{i,j}} = \ker L_{\phi_{i,j}}\quad\text{(modulo infinitesimal symmetries of \(S^2\) and target isometries)}.
\]
Let 
\[
\nu_i := \lim_{r\to0}\lim_{k\to\infty}\Big(E(u_k;A_{r,\rho}(x_i)) - \sum_j E(\phi_{i,j})\Big)
\]
be the neck energy at \(x_i\).

**Conjecture (Jacobi–neck correspondence):**  
If every bubble \(\phi_{i,j}\) appearing in the blow-up analysis is *nondegenerate* (that is, \(\mathcal K_{\phi_{i,j}} = \{0\}\)), then one has
\[
\nu_i = 0 \quad \text{for every blow-up point } x_i,
\]
for all weakly controlled approximate harmonic sequences \(u_k\) with \(f_k \to 0\) in \(H^{-1}\).

Equivalently: nondegeneracy of all harmonic spheres forbids the appearance of nonzero neck defect measures in any weak limit of approximate harmonic maps from surfaces.

More generally, if some \(\phi_{i,j}\) has nontrivial \(\mathcal K_{\phi_{i,j}}\), can one produce, via an explicit perturbative construction along a slowly varying tangent direction in \(\mathcal K_{\phi_{i,j}}\), approximate harmonic maps with a prescribed small nonzero neck energy \(\nu_i>0\)?  

The conjecture seeks to identify the presence (or absence) of kernel modes of Jacobi operators as the *sole geometric mechanism* responsible for residual neck energy in bubbling sequences.

Why is it a "good" problem:  
This refined conjecture isolates a precise equivalence between an analytic compactness property (energy quantization without residual neck energy) and a geometric–spectral property (nondegeneracy of all bubble spheres).  It is simple to state yet lies far beyond current analytical techniques: although energy identities are known under certain curvature or integrability conditions, no result connects them to the spectral structure of the linearized operator on the bubbles.  
Establishing (or disproving) this Jacobi–neck correspondence would create a conceptual bridge between two deep areas—spectral geometry of harmonic spheres and defect‑measure analysis in geometric variational problems—yielding a potential unification of linearized stability theory and nonlinear concentration phenomena.  The problem is mathematically sound, sharply formulated, and potentially transformative, capturing the essence of a “good” research-level mathematical problem.